\cotefile{14}
\scanpage{59}
est également caractérisé par la condition que le diagramme suivant soit commutatif

\diagram{\begin{tikzcd}[ampersand replacement=\&,column sep=large,row sep=large]
H^1(A_Y) \arrow[r,"u^{(1)}"] \arrow[d,"\varphi_Y^{\unsure{(1)}}"',"{\substack{\text{isom mod}\\\text{torsion}}}"] \& H^1(B_X) \\
H^1(Y) \arrow[r,swap,"{a(\gamma^{\deleted{\unsure{r}}\unsure{n+m-1|m+m-1}}(z))}"] \& H^{2n-1}(X)(n-1) \arrow[u,swap,"\psi_X"]
\end{tikzcd}}

(où la cohomologie est à coefficients dans $\mathrm{Z}_l$, $l$ premier à la caractéristique). Enfin, $u$ ne dépend que de la classe d'équivalence \uline{numérique} de $\deleted{\unsure{z|\mathfrak{z}}}Z$.

Démonstration. Considérons l'homomorphisme composé $\psi_X\,a(\gamma(z))\,\varphi_X^{(1)}$. Il provient d'une classe de \deleted{\unsure{équival|équivalence}} correspondances algébriques (indépendante de $l$), car les trois facteurs sont de ce type. En vertu de 1.7. il est de la forme $u^{(1)}$, où $u:B_X\longrightarrow A_Y$ est un hom de V.A. Il reste à prouver la formule
\[
u(\operatorname{cl}(D))=S_Y({}^tZ(D))
\]
lorsque $ {}^tZ(D)$ est défini (le fait que cette formule caractérise $u$ est immédiat, car pour tout $\zeta\in\operatorname{Pic}^{\unsure{0}}(X)$, il y a un représentant $D$ tel que $ {}^tZ(D)$ soit défini). On va
% archive sheet number: 58

\scanpage{60}
d'abord construire un $u:B_X\longrightarrow A_Y$ ayant cette dernière vertu, indépendamment de considérations cohomologiques. \deleted{Soit un diviseur sur $X\times B_X$}\inserted{Notons d'abord qu'il suffit de} faire le calcul avec les classes d'équivalence linéaire $z$, $d$ des cycles (ou $\unsure{\overline{\mathrm m}|=}$ l'équivalence d'Albanese) -- puisque si $\unsure{W|\Omega|\mathfrak{A}}\in\operatorname{Cyc}_0(Y)$ est lin. équi. à $0$, $S_Y(W)$ est nul. Donc on n'a pas besoin de se préoccuper de la question \deleted{que les produits}\inserted{de \deleted{\unsure{de}} l'existence} des \unsure{expressions}, et se borner à regarder $d\longmapsto S_Y(z(d))$. Or le calcul des intersections prouve que, si $\pi:Y\times X\times B_X\longrightarrow Y\times B_X$ est la projection canonique, \inserted{et $\Theta\in A^1(X\times B_X)$ l'élément canonique,} on a \deleted{(identifiant $d$ à un point}
\[
z(\Theta(b))\simeq\deleted{\unsure{\pi_*}}C(b)
\]
où $C\in Z^m(Y\times B_X)$ est donné par
\[
C=\pi_*\bigl((1_Y\times\Theta)\mathbin{.}(z\times1_{B_X})\bigr)
\]
\inserted{D'ailleurs $z(\Theta(0))=0$ donc $z(\Theta(b))=z(\Theta((b)-(0)))$}

On applique alors 1.1. à $C$ sur $Y\times B_X$, pour obtenir $u:A_{B_X}=B_X\longrightarrow A_Y$, satisfaisant
\[
\begin{aligned}
&\deleted{\unsure{u\,z(\Theta(b))|u\,\mathfrak{z}(\Theta(b))}}\\
&S_Y(C((b)-(0)))=u(b)\qquad\text{i.e.}\\
&S_Y(z(\Theta(b)))=u(b).
\end{aligned}
\]
% archive sheet number: 59

\scanpage{61}
Le fait que cela donne ce qu'il faut du point de vue cohomologique se déduit par exemple de 1.6. en suivant la construction précédente du point de vue cohomologique. Bien entendu, $u$ dépend linéairement de $z$, et si $z$ est num. \unsure{équi.|équiv.} à $0$, il en est de $\unsure{\overline{\mathrm m}|=}$ de $C$, donc de $u$ en vertu de 1.1.

Variante de la démonstration. On cherche $B_X\unsure{\mathord{-}|\to} A_Y\simeq B_Y^*$, i.e. une \unsure{clas.|classe|cl.} de corr. div. sur $B_X\times B_Y$. On la construit en termes de $\Theta_X\in A^1(X\times B_X)$, $\Theta_Y\in A^1(Y\times B_Y)$, $Z\in A^{n+m-1}(X\times Y)$, \deleted{par le procédé en \unsure{y} considérant $\unsure{\pi}:X\times B_X\times\inserted{\Theta_X}\unsure{Y}\times B_Y$ et} en formant la classe de \unsure{correspondance|correspondances} $\Theta_Y\circ Z\circ{}^t\Theta_X$ sur $B_X\times B_Y$,
\[
\ldots\,-
\]

\deleted{\uline{Corollaire} \inserted{2.2} Tout homomorphisme $B_X\mathord{-}A_Y$ est défini par un $Z_1\in\operatorname{Cyc}_1(X\times Y)$.}
% archive sheet number: 60

\scanpage{62}
\uline{Corollaire} 2.\unsure{?|2}. \margin{\rule[-2.8\baselineskip]{0.45pt}{3.5\baselineskip}}Notations de 2.1. Considérons $ {}^t(u_Z):B_Y\longrightarrow A_X$, alors $ {}^t(u_Z)=u_{({}^tZ)}$, i.e. on a pour tout diviseur $D$ alg. équiv. à $0$ sur $Y$, tel que $Z(D)$ soit défini,
\[
S_X(Z(D))={}^t(u_Z)(D)
\]
% slanted
\margin{Attention aux signes éventuels !}
Vérification par exemple par voie cohomologique.

$[$ \uline{Proposition} 2.3. Conditions de 1.1. \deleted{Alors on a}\inserted{Considérons} $v_Z={}^t(u_Z):B_X\longrightarrow B_Y$. Alors on a, pour tout diviseur alg. équiv. à $0$ $D$ sur $X$, tel que $Z(D)$ soit défini,
% slanted
\margin{devrait remonter au N\textsuperscript{o}1.}
\[
{}^tZ(D)={}^t(\unsure{u|v}_Z)(D)
\]

N.B. Ceci doit pouvoir servir à une autre démonstration de 1.1., via un hom de schémas en groupes $\underline{\operatorname{Pic}}_{X/k}\longrightarrow\underline{\operatorname{Pic}}_{Y/k}$ !
% slanted
\margin{attention aux signes éventuels !}

Vérification de 2.3. par AQT $]$

\uline{Proposition} 2.4. Avec les notations de 2.1. \deleted{Supposons que}\inserted{Fixons $l$,} pour \deleted{un $l$ convenable,}\inserted{que} l'hom. $a(\gamma(z)):\allowbreak H^1(Y,\mathbb Q_l)\longrightarrow H^{2n-1}(X,\mathbb Q_l)(n-1)$ soit un isom \deleted{\unsure{modulo|mod.} torsion}\inserted{il faut et s. que $u_Z$ soit une isogénie (cond. ind. de $l$)} \deleted{\unsure{diviseurs|et suffisante}}. Alors il existe un \deleted{cycle}diviseur
% archive sheet number: 61

\scanpage{63}
\deleted{algébrique}\inserted{$D$} à coefficients \uline{rationnels}, \deleted{\unsure{dans}}\inserted{sur} $X\times Y$ tel que l'hom. $\deleted{\unsure{a(z)}}a({}^t\gamma(D)):H^{2n-1}(X,\mathbb Q_l)(n-1)\longrightarrow H^1(Y,Q_l)$ soit inverse de $a(\gamma(z))$, ce qui caractérise $D$ comme élément de $\unsure{\operatorname{Cl}(X\times Y)|\operatorname{Cl}X\times Y)|\mathcal C\ell(X\times Y)}\otimes Q$. L'\deleted{hom}élément de $\operatorname{Hom}_{Q}(A_Y,B_X)$ défini par $D$ n'est autre que $u_Z^{-1}$, \deleted{\unsure{à un signe}} Pour qu'on puisse choisir $D$ à coefficients entiers, il faut et il suffit que $u$ soit un isom. \deleted{Pour qu'on} Plus généralement, si $\operatorname{Ker}u$ est de cardinal $N$, \deleted{\unsure{on p}} (Attention, $\operatorname{Ker}u$ en tant que \uline{schéma} en groupes) alors on peut choisir \inserted{$D$} de la forme $D_{\unsure{0}}=\deleted{\unsure{p}}N^{-r}D_0$, $D_0$ étant un div. à coeff. entier, et $r\geq0$, \deleted{\unsure{sss} l'entier $m\geq1$}\inserted{mais pas sous la} forme $\deleted{\unsure{p}}M^{-r}D_0$, \deleted{\unsure{dès|doit}} $\unsure{m|M}$ ne contient pas tous les diviseurs premiers de $n$.

On écrit le diagramme comm
\diagram{\begin{tikzcd}[ampersand replacement=\&,column sep=large,row sep=large]
H^1(Y,Q_l)\arrow[r,"{a(\gamma(z))}"] \& H^{2n-1}(X,Q_l)(n-1)\arrow[d,"\unsure{\wr|2}"] \\
H^1(A_Y,Q_l)\arrow[u,swap,"\wr"]\arrow[r,"u^{(1)}"] \& H^1(B_X,Q_l)
\end{tikzcd}}
où \unsure{toutes} les flèches écrites sont $\mathbb Z$-algébriques, les flèches verticales des isom. \deleted{\unsure{Si}} Pour que
% archive sheet number: 62

\scanpage{64}
$a(\gamma(z))$ soit un isom, il f et s que $u^{(1)}$ le soit, i.e. que $u$ soit une isogénie. Alors $u$ admet un inverse $u^{-1}$ unique dans $\operatorname{Hom}_Q(A_Y,B_X)$, \deleted{défini} donc l'hom inverse de $u^{(1)}$ est défini par un $Q$-cycle algébrique. \inserted{$R$, le meilleur dénominateur possible \inserted{étant donné par $\operatorname{card}\operatorname{Ker}u$}}\deleted{\unsure{coeff. en|à coeff. en}} Donc \unsure{$(a(\gamma(z))^{-1})$|$(a(\gamma(z))^{-1}$|$a(\gamma(Z))^{-1}$} est défini par un $Q$-cycle algébrique, ce qui prouve la première assertion, \deleted{La seconde, se voit comme}\inserted{ayant $\unsure{\overline{\mathrm m}|=}$} dénominateur que $R$. Cela prouve la première assertion, la seconde résulte de la connaissance de l'effet cohomologique d'une classe de \inserted{corr.} \unsure{divisoralles|divisorielles}, via l'effet sur les VA \ldots{} (cf \ldots{} ). Il reste à \deleted{prouver}\inserted{noter} que si $D$ admet un certain dénominateur, ce dernier vaut pour le \deleted{hom} $Q$-hom $A_Y\unsure{\mathbin{-}|\to}B_X$ inverse de $u$, donc il \deleted{\unsure{ne cont}} ne peut être meilleur qu'une puissance de \deleted{card} $N$.

\uline{Corollaire} 2.5. Avec les notations de 2.4., pour
% archive sheet number: 63

\scanpage{65}
que \inserted{$u$ soit une isogénie et que} $N$ soit premier à $l$ (où $l\ne\operatorname{car}k$) il f. et s. que le composé
\[
H^1(Y,\mathbb Z_l)\unsure{\mathord{-}|\longrightarrow} H^{2n-1}(X,\mathbb Z_l)(n-1)\mathord{-}
\frac{H^{2n-1}(X,\mathbb Z_l)(n-1)}{\text{torsion}}
\]
soit \deleted{surjectif} bijectif; \deleted{et}\inserted{Pour que de plus $H^{2n-1}(X,\mathbb Z_l)$ soit sans torsion,} il suffit pour ceci que $H^1(Y,\mathbb Z_l)\to H^{2n-1}(X,\mathbb Z_l)(n-1)$ soit bijectif.

% slanted
\margin{\uline{Th} ? C'est un cas particulier important du ``Lefschetz \unsure{vache|vachement}'' \unsure{donnant|dérivant} Lefschetz des surfaces\ldots.}
\uline{Corollaire} 2.6. Soit $X$ projective lisse conn. de \unsure{dim} $n\geq2$, et $\xi_1,\ldots,\xi_{n-1}\in\inserted{\operatorname{Pic}(X)}$ \deleted{\unsure{classes des cat. ly lin} \unsure{polarisat|provenant}} \deleted{\unsure{une telle}} des faisceaux amples, \inserted{$\zeta_1,\ldots,\zeta_{n-1}\in A^1(X)$ leurs classes.} Alors pour tout entier $l$, l'hom de cup-produit par
\[
\deleted{\unsure{\xi}}\eta=\zeta_1\zeta_2\ldots\zeta_{n-1}:
\]
\[
H^1(X,\deleted{\unsure{Z}}Q_l)\longrightarrow H^{2n-1}(X,Q_l)(n-1)
\]
est un \uline{isomorphisme}. L'isom inverse $L$ est définissable par un élément de $\deleted{\unsure{Z^1_{\mathrm{rat}}}}A^1(X\times X)\otimes Q$, \deleted{i.e. de} ou \unsure{même|encore} de $\operatorname{Cl}(X,X)\otimes Q$, et \deleted{cet élément} \inserted{\unsure{ce|cet}} dernier élément est unique et ne dépend pas de $l$.

Il faut prouver que l'hom de VA
\[
u:B_X\longrightarrow A_X
\]
% archive sheet number: 64

\scanpage{66}
``correspondant'' est une isogénie, ou encore, a un noyau fini. Pour bien faire, \deleted{il faudrait}\inserted{on peut} prouver mieux, savoir qu'il correspond à une \uline{polarisation} de $B_X$ (compte tenu de $A_X\simeq B_X^*$).

\deleted{Or cet hom. peut se décrire de proche en proche comme composé d'hom \deleted{de restriction},}

\diagram{\begin{tikzcd}[ampersand replacement=\&,column sep=tiny,row sep=small,execute at end picture={\draw[thin] (current bounding box.south west) -- (current bounding box.north east);}]
\deleted{B_X=B_{X_0}}\arrow[r,"{\deleted{\rho_0}}"] \& \deleted{B_{X_1}}\arrow[r,"{\deleted{\rho_1}}"] \& \deleted{B_{X_2}} \& \deleted{\ldots}\arrow[r,"{\deleted{\rho_{n-3}}}"] \& \deleted{B_{X_{n-2}}}\arrow[r,"{\deleted{\sigma}}"] \& \deleted{A_{X_{\unsure{n-2}}}}\arrow[r,"{\deleted{i}}"] \& \deleted{\unsure{\vert|1}}\\
\& \& \& \& {} \arrow[r] \& \deleted{B_C\simeq A_C}\arrow[r] \& {}
\end{tikzcd}}
\diagram{\begin{tikzcd}[ampersand replacement=\&,execute at end picture={\draw[thin] (current bounding box.north west) -- (current bounding box.south east);}]
\deleted{B_X}\arrow[r,"{\deleted{\rho}}"] \& \deleted{B_Y}\arrow[r,"{\deleted{\lambda}}"] \&\deleted{A_Y}\arrow[r,"{\deleted{\sigma}}"] \&\deleted{A_X}
\end{tikzcd}}

\deleted{où $Y$ est défini \deleted{par} comme l'intersection de $n-2$ \deleted{\unsure{hypersup|hyperplans}} diviseurs ``généraux'' correspondant à $L_1,\ldots,L_{n-2}$ (supposés \uline{très} amples, ce qui est loisible) $\rho$ est la restriction, $\sigma$ l'injection, et $\lambda$ est défini comme $u$, en termes de $L_{n-1}\mid Y$ sur $Y$. On sait \inserted{(de proche en proche\ldots)} que $\rho,\sigma$ sont des isogénies, \deleted{et} (d'ailleurs transposées \deleted{\unsure{sg}} l'une de l'autre),}

Or une fonctorialité évidente en $X,Y$ permet d'interpréter de la façon suivante l'hom $B_X\unsure{\mathord{-}|\to}A_Y$ défini dans 2.1., lorsque $Z$ est défini par un morphisme
% archive sheet number: 65

\scanpage{67}
\[
\unsure{v|\nu}: C\longrightarrow X\times Y,
\]
où $C$ est une courbe lisse projective, donc par
\[
v_1:C\unsure{\mathbin{-}|\to}X,\qquad v_2:C\mathbin{-}Y
\]
Il suffit alors de prendre le composé
\diagram{\begin{tikzcd}[ampersand replacement=\&]
B_X\arrow[r,dash,"v_1^*"] \& B_C\arrow[r,"\text{dualité}"] \& A_C\arrow[r,"v_{2*}"] \& A_Y
\end{tikzcd}}
ou encore, en termes de classes de correspondances divisorielles sur $B_X\times B_Y$, il suffit de composer la classe \deleted{\unsure{du diviseur|dualisante}} canonique sur $B_C\times B_C$, avec les hom $B_X\xrightarrow{v_1^*}B_C$, $B_Y\overset{v_2^*}{\mathbin{-}}B_C$. Ceci posé, si $X=Y$, $v_1=v_2$, on voit que \deleted{cette classe de} les conditions suivantes sont équivalentes :

(i) L'hom $B_X\mathbin{-}A_X=B_X^*$ est un hom de polarisation

(ii) L'hom $B_X\underset{\unsure{\cdot|\prime}}{\mathbin{-}}B_X^*=A_X$ est une isogénie

(iii) L'hom $\unsure{v|w}^*:B_X\mathbin{-}B_C$ est injectif \inserted{mod groupes finis.}

(iv) L'hom $v_*:A_C\mathbin{-}A_X$ est surjectif
% archive sheet number: 66

\scanpage{68}
D'ailleurs, si $\unsure{c|\tau|\xi}\in\deleted{\unsure{v}}A_1(X)$ est défini par $\unsure{u|v}:C\unsure{\mathbin{-}|\to}X$, alors $B_X\to A_X$ est aussi, au pt de vue cohomologique, défini par le cup-produit par
\[
\begin{aligned}
&\gamma(c):H^1(A_X,Q_l)\xrightarrow{\sim}H^1(X,Q_l)\\
&\longrightarrow H^1(B_X,Q_l)\xleftarrow{\sim}H^{2n-1}(X,Q_l)(n-1).
\end{aligned}
\]
$/$ On est ramené au fait bien connu que si $C$ est une courbe sur $X$, intersection ``gén.'' de diviseurs très amples, alors $A_C\mathbin{-}A_X$ est surjectif

\uline{Proposition} 2.7. $X$ et $Y$ comme d'habitude, $D\in\deleted{Z}A^1(\unsure{x|X}\times Y)$, d'où
\[
u_D:A_Y\longrightarrow\deleted{\unsure{A_C}}B_X
\]
et
\[
a(\gamma(D)):H^{\deleted{\unsure{1}}\inserted{2n-1}}(X,\mathbb Q_l)(n-1)\mathbin{-}H^1(Y,\mathbb Q_l)
\]
Pour que $a(\gamma(D))$ soit un isom, il f. et s. que $u_D$ soit une isogénie. Alors il existe $Z\in\deleted{Z}A_1(x\times Y)\otimes Q$ tel que $\deleted{\unsure{a_z|{}^t}}a(\gamma(z)):H^1(Y,Q_l)\longrightarrow H^{\deleted{\unsure{1}}\inserted{2n-1}}(X,Q_l)(n-1)$ soit \deleted{défini} inverse de $a(\gamma(D))$ (quel que soit $l$). \deleted{Ce $Z$ est}
% archive sheet number: 67

\scanpage{69}
\deleted{\unsure{Le $Q$-hom|L'élément|Rappelons}} le $Q$-hom \deleted{On écrit le diagramme \unsure{com}}
\[
u_Z:B_X\longrightarrow A_Y\qquad(\text{mod isog})
\]
défini par $Z$ en vertu de 2.1. n'est autre que $u_D^{-1}$.

\uline{Dém}. On utilise le diagramme commutatif
\diagram{\begin{tikzcd}[ampersand replacement=\&,column sep=large,row sep=large]
H^{2n-1}(X,Q_l)(n-1)\arrow[r,"{a(\gamma(D))}"]\arrow[d,"\unsure{\wr|2}\,\deleted{\unsure{\text{isom}}}","\psi_X"'] \& H^1(Y,Q_l)\\
H^1(B_X,\mathbb Q_l)\arrow[r,"u_D^{(1)}"] \& H^1(A_Y,Q_l)\arrow[u,swap,"\wr\quad\varphi_Y"]
\end{tikzcd}}
\unsure{qui|Qui|Ceci} prouve que $a(\gamma(D))$ est un isom. ssi $u_D^{(1)}$ l'est, i.e. $u_D$ est une isogénie, Pour conclure à l'existence de $Z$, il faut montrer que $\varphi_Y^{-1}$ et $\psi_X^{-1}$ sont $Q$-algébriques (de façon indépendante de $\unsure{l|k}$). \deleted{Pour ceci, il faut \unsure{pr}}\deleted{\inserted{On le sait déjà}} Pour $\psi_X$, on est donc ramené à prouver 2.7. dans le cas particulier de l'hom can.
\[
H^{2n-1}(X,Q_l)(n-1)\longrightarrow H^1(B_X,Q_l)
\]
défini par $D\in\operatorname{Cl}(X,B_X)$, pour trouver un élément can. $Z\in A_1(X\times B_X)$.
% archive sheet number: 68

\scanpage{70}
Or choisissons un $\zeta\in A^1(X)$ de polarisation, Alors on a un diagramme comm.
\diagram{\begin{tikzcd}[ampersand replacement=\&,column sep=large,row sep=large]
H^{2n-1}(X,Q_l)(n-1)\arrow[r,"{\substack{\psi_X\\\sim}}"] \& H^1(B_X,Q_l)\\
H^1(X,Q_l)\arrow[u,"{\lambda\quad\unsure{\wr|2}}","{\text{cup }\zeta^{n-1}}"'] \& H^1(A_X,Q_l)\arrow[l,"\varphi_X^{\deleted{\unsure{(1)}}}","\sim"']\arrow[l,phantom,"{\deleted{\rightarrow}}"{description,pos=.08}]\arrow[u,swap,"{\wr\quad\mu}"]
\end{tikzcd}}
où toutes les flèches écrites \deleted{\unsure{sont}} \deleted{\unsure{et}} définies par des classes de corr. alg. (ind. de $l$). \inserted{($\mu$ étant défini comme composé \inserted{$\psi_X\lambda\varphi_X$}) et sont des isom. en vertu de \ldots{} et de 2.6.} Donc $\psi_X^{-1}=\lambda\varphi\mu^{-1}$, et tout revient à voir que $\mu^{-1}$ est algébrique, ce qui provient du fait que $\mu$ est défini par un $Q$-\unsure{isom|hom|isogénie} $B_X\unsure{\mathbin{-}|\to}A_X$ \inserted{(cf}, dont on prend simplement l'inverse\ldots{}

Reste à traiter le cas de $\varphi_Y^{-1}$ (où on peut remplacer $X$ par $Y$). On utilise le $\unsure{\overline{\mathrm m}|=}$ diagramme, qui donne
\[
\varphi_X^{-1}=\mu^{-1}\psi_X\lambda
\]
donc on est ramené à inverser $\mu$ \unsure{``algébriquement''|``algébriquemt''|``algébriques''}, or cela est possible grâce au fait que $\mu$ est associé à une isogénie de V.A.
% archive sheet number: 69

\scanpage{71}
\uline{Questions} 1) Y-a-t-il tjrs un cycle \uline{entier} qui donne l'inverse pour $H^1(A_X)\to H^1(X)$ ? Il suffirait qu'on sache que pour $n$ \deleted{grand} convenable (\unsure{donc|alors} $n$ grand) la fibre générique de $\operatorname{Sym}^n(X)\to A_X$ a un point rationnel.

2) \unsure{Y-at-il|Y-a-t-il} tjrs un cycle \uline{entier} qui donne l'inverse pour $H^{2n-1}(X)(n-1)\to H^1(B_X)$ ? \deleted{\unsure{Il}}Si la réponse à 1) était affirmative, il suffirait de connaître la réponse à 2) dans le cas des VA.
% archive sheet number: 70

\scanpage{73}
\[
H^{\unsure{1|1'}}(X,\mu_N)\simeq {}_N\operatorname{Pic}(X)
\]
\[
0\to {}_N\operatorname{Pic}^0(X)\to H^1(X,\mu_N)\to {}_N\mathrm{NS}(X)\to 0
\]
\margin{\tikz[baseline=-.4ex]{\draw[->] (.22,-.2) -- (0,-.2) -- (0,.85);\draw (0,.3) -- (.28,.3);} $X$ propre sur corps $k$ sép. clos, $N,l$ premiers à car. de $k$}

\[
H^1(X,\mathbb Z_l(1))\xleftarrow{\sim}
\deleted{H^1(X,A_X)}T_l(\underline{\operatorname{Pic}}{}^0_{X/k})\deleted{\unsure{H^1|B}}
\]
[Utiliser le fait que $\operatorname{Pic}^{\tau}(X)/\operatorname{Pic}^0(X)\deleted{\underline{\operatorname{Pic}}_{X/k}}$ est \uline{fini}]

En particulier, si $\underline{\operatorname{Pic}}{}^0_{X/k}$ propre, p.ex. $X$ normal, on trouve, introduisant $B_X=\underline{\operatorname{Pic}}{}^{00}_{X/k}$ et $A_X=B_X^*$, on a
\[
H^1(X,\mathbb Z_l\deleted{\unsure{(1)}})\xleftarrow{\varphi_X^*}H^1(A_X,\mathbb Z_l).
\]

\deleted{$H^{2n-1}$} Supposons $X$ lisse, purement $n$-dim. alors
\[
H^{2n-1}(X,\mathbb Z/N\mathbb Z)\simeq H^1(X,\mathbb Z/N\mathbb Z)^\vee(-n)
\]
d'où
\diagram{\begin{tikzcd}[ampersand replacement=\&,column sep=small]
0\arrow[r]\& \deleted{\unsure{B|P}}({}_N\mathrm{NS}(X))^\vee\deleted{\unsure{(-n)|(+n)}}\arrow[r]\&
 H^{2n-1}(X,\mathbb Z/N\mathbb Z)\arrow[r]\& {}_N A_X(k)(-n)\arrow[r]\arrow[d,dash,"\sim"']\&0\\
\&\&\& H^1(B_X,\mathbb Z/N\mathbb Z)\deleted{\unsure{(-n)|(+n)}}\inserted{(1-n)}\&
\end{tikzcd}}

Passant à la limite sur $l^{\unsure{\nu|n}}$
\diagram{\begin{tikzcd}[ampersand replacement=\&,column sep=small,row sep=small]
0\arrow[r]\&( \mathrm{NS}(X)[l])^\vee\deleted{\unsure{(-n)|(n)}}\arrow[r]\arrow[d,"\phantom{l}"]\&
H^{2n-1}(X,\mathbb Z_l)(n-1)\arrow[r,"\text{hom algébrique}"']\&
H^1(B_X,\mathbb Z_l)\deleted{\unsure{(1-n)|(-n)}}\arrow[r,dash]\&0\\
\&\substack{\text{comp.}\\ l\text{-primaire}}\&\&\&
\end{tikzcd}}
% archive sheet number: 72

